\section{Ontologies and Knowledge Representation}
\begin{itemize}
    \item Definition of an ontology: concepts, relations, and formal structure
    \item Motivations for using ontologies: shared vocabulary, explicit assumptions, reusability
    \item Relevance for ML method selection: structuring hierarchical relationships between 
    methods, tasks, and problem types to support interpretable recommendations
    \item Ontologies in the ML domain: existing ML ontologies and their properties
    \item Limitations: manual effort, iterative maintenance, scope constraints
\end{itemize}

An ontology is an explicit specification of a conceptualization,
where a conceptualization is an abstract model of a domain that
identifies its relevant concepts and the relationships between them
\parencite{gruber1995toward_principles_design_ontologies_used_knowledge_sharing}.
\textcite{studer1998knowledge_engineering_principles_methods} refine
this by describing an ontology as a
\textit{formal, explicit specification of a shared conceptualization}.
Each term in this definition carries a distinct meaning.
\textit{Formal} means the ontology is machine-readable and expressed
in a language with precise semantics.
\textit{Explicit} means that all concepts and the constraints on
their use are clearly defined, rather than left implicit in code or
documentation.
\textit{Shared} means the ontology captures knowledge that is
accepted by a group, rather than reflecting the perspective of
a single individual
\parencite{studer1998knowledge_engineering_principles_methods}.
In practical terms, an ontology defines a common vocabulary for a
domain, including machine-interpretable definitions of its basic
concepts and the relations among them
\parencite{noy2001ontology_development_101_guide_creating_first_ontology}.

An ontology consists of classes, which represent concepts in the
domain, properties that describe features of those concepts, and
relations that link concepts to one another
\parencite{noy2001ontology_development_101_guide_creating_first_ontology}.
Classes are typically organized in a hierarchy where subclasses
inherit the properties of their parent classes.
Taken together, these elements allow a domain to be described in a
way that both humans and software can interpret and reason over
\parencite{gruber1995toward_principles_design_ontologies_used_knowledge_sharing}.

The design of an ontology involves several practical constraints.
\textcite{gruber1995toward_principles_design_ontologies_used_knowledge_sharing}
argues that an ontology should require only the minimal ontological
commitment needed for its intended purpose, meaning it should make
as few claims about the world as necessary.
This principle reflects a trade-off identified by
\textcite{studer1998knowledge_engineering_principles_methods}:
the more reusable an ontology is, the less directly applicable it
tends to be in practice, since a more general representation
provides less operational guidance for a specific task.
\textcite{noy2001ontology_development_101_guide_creating_first_ontology}
point to a related constraint on scope.
An ontology should not attempt to capture everything that could be
said about a domain, but only what is needed for the application at
hand.
Expanding scope beyond what the application requires adds complexity
without benefit, and may introduce inconsistencies that are difficult
to resolve.
Ontology development is also not a one-time activity.
As domain knowledge evolves and new requirements emerge, the
ontology must be revised, which means the initial design should
anticipate extension and allow the vocabulary to be specialized
without requiring changes to existing definitions
\parencite{gruber1995toward_principles_design_ontologies_used_knowledge_sharing,
noy2001ontology_development_101_guide_creating_first_ontology}.